\section{The Best Possible World}

\begin{quotex}
If we assume the existence of an omniscient and omnipotent being, one that knows and can do absolutely everything, then to my own very limited self, it would seem that existence for it would be unbearable. Nothing to wonder about? Nothing to ponder over? Nothing to discover? Eternity in such a heaven would surely be indistinguishable from hell.
\flright{\textsc{Isaac Asimov}, \emph{``X'' Stands for Unknown, 1984}}
\end{quotex}


Apart from his embarrassingly juvenile understanding of God, Asimov raises a good question, viz.: What is the best possible world? Such a world would be wonderous, and offer mysteries to discover. Ultimately, we would like to discover the most wonderous thing of all, but let's hold that off for a moment.

Asimov's point differs from others like him. The most common argument against God is that this world, with all its horrors and evils, could not have been created by an omniscient, omnipotent, and good God. You will note that Asimov does omit the qualifier ``good''.

Let's consider the creation of the world as described in Genesis. Note that it is more like a framework rather than a detailed description. Instead of 100, or many more, volumes with tedious details, Genesis includes just a few short chapters. That leaves plenty to discover about the creation of the world.

Spiritual texts assume a rather Platonic view of the world, provided that is understood in a very general sense. Plato is just one link of a metaphysical chain that forms the Heart of the West\footnote{\url{https://gornahoor.net/?p=1256}}. The fundamental point is that, in addition to the material world, there is a spiritual world of ideas. A ``thing'', therefore, is a combination of the physical with an idea. There is no definition in physics for any particular thing. That is called the ``measurement problem'' and is still rather mysterious. Quantum physics describes a state vector for particles that can do no more than describe a probability. There is no physical formula for a star, a plant, or a human being.

\paragraph{The Stages of Genesis}


Genesis describes several stages of creation which match our personal experiences perfectly.

\begin{itemize}


\item \textbf{Possibilities}. The world is divided into Heaven (Ideas) and Earth (matter).

\item \textbf{Movement}. The world is set in motion. By Newton's First Law, matter cannot begin to move without an outside force acting on it.

\item \textbf{Vegetation}. Life begins.


\item \textbf{Animals}. Sentient life begins.


\item \textbf{Humans}. Rational life begins.


\item \textbf{Freedom}. Humans are given the gift of free will.


\item \textbf{Love}. Only free beings can love

\end{itemize}


In your perfect world, what would you change? To be sentient means to be able to ``feel''. An animal feels thirsty, so it seeks water. A robot, running low on batteries, just reads a gauge, which is not really part of him.

Asimov's world requires rational life, beings who are able to discover things. If we knew everything, as he says, there is nothing to discover. Since we don't, we have to stumble through the world with imperfect knowledge. The price of imperfect knowledge includes mistakes, often costly mistakes. But it also includes creativity.

Would the best world deny freedom? Perhaps we could have been programmed like robots, with something like Asimov's Three Laws of Robotics\footnote{\url{https://en.wikipedia.org/wiki/Three\_Laws\_of\_Robotics}}. Then we would always make the moral choice. The world envisioned by Dostoevsky's Grand Inquisitor\footnote{\url{https://www.mtholyoke.edu/acad/intrel/pol116/grand.htm}} provided for man's needs while denying him freedom.

Finally, there is Love. Would a perfect world include the ability to Love? Only free beings can love. Many people today, in the anonymous loneliness of the modern world, substitute for it the affection of dogs, which cannot love in any satisfactory sense.

Love encompasses all the aspects of Being:

\begin{itemize}


\item There are intense sexual sensations at the most primitive level of life.

\item The is the exhilarating feeling of falling in love with someone.

\item There is the feeling associated with creativity such as that of an artist, mathematician or scientist.

\item There is the thrill at understanding the loftiest philosophical understanding as described, for example, by Spinoza's intellectual love of God.

\item Ultimately there is the love of God with all one's heart, soul, strength, and mind.
\end{itemize}


I suppose most everyone has a gripe against God and wishes things could be different. But a world with life, feelings, discovery, creativity, freedom, and love must necessarily be better than a computerized simulation of life or a world populated by programmed robots.

\url{https://www.youtube.com/watch?v=Tdx6lLvvRyg} 

\flrightit{Posted on 2022-05-22 by Cologero}

\begin{center}* * *\end{center}

\begin{footnotesize}
\begin{sffamily}
\texttt{Dochxa Rhoan on 2022-05-23 at 13:45 said:}

My first comment here though I've been following for several months. Thanks for many thought-provoking discourses.

Regarding ``free will'' -- been reading some neurological experiments and results seem to increasingly indicate a distinction required between predetermined actions arising out of a particular neurological state and the experience of free will. (This underscores, I think, the material world/idea world connection or absence of connection.)

I'm not arguing for an assertion or denial of free will; rather, for clarification of exactly what is being named when we say ``humans have free will''.

Thanks in advance.

\hfill

\texttt{Cologero on 2022-05-23 at 14:22 said:}

You can start here for the metaphysical understanding:

\textit{Principial Origin and Final Destiny}\footnote{\url{https://www.gornahoor.net/?p=15342}}

\textit{Buridan's Donkey}\footnote{\url{https://www.gornahoor.net/?p=15031}}

\textit{Did you ever have to make up your mind?}\footnote{\url{https://www.youtube.com/watch?v=C0cJ__JBnWg}}

Then type ``free will'' into the little search box

In your spare time, try this experiment:

Take 12 subjects and promise them a free ice cream for participating

Hook them up to some electrodes and let the experimenter determine which flavor ice cream each subject will choose

Let me know the success rate

There is no ``experience of free will''. That little conscious part of yourself is NOT your true self.

I plan to start a business and am looking for suckers investors.

It is a home free will detector kit. You hook yourself up in the morning and you will know everything you will do for the rest of the day!

Send donations to my patreon account and remember that past performance is no indicator of future results.

\hfill

\texttt{Dochxa Rhoan on 2022-06-01 at 22:23 said:}

``There is no experience of free will'' -- which is why I asked for some definition. That little conscious part of myself asserts that I choose things, but, stepping out of my little conscious part of myself frame of reference, I suspect I do not. The ass experiment is a good one -- what about the flipping of a coin, or an electron passing through a slot? Isn't it the same thing except our little conscious part of ourselves doesn't have the full equation to determine the outcome? I'm not making a challenge, mind you, except that jumping back and forth from omniscience hardly seems the best way to forward to discourse. I show up in the world in a particular way regardless (to what's call ``me'') of the unconscious mechanisms which may be predetermined. From the point of view of God, to ``me'' it is moot. From my mind's point of view, I know it is futile. In my being present to the world, however, it seems as though there is something to explore. Cheers!

\hfill

\texttt{Cologero on 2022-06-01 at 23:31 said:}

Fascinating, Dochxa, but not even close. I absolutely did provide a link to a precise metaphysical understanding of free will. You cannot reason your way to free will, you have to will it. People seem to enjoy running around in circles like hamsters on a wheel.

It is true, though, that free will has to be earned.

\hfill

\texttt{Dochxa Rhoan on 2022-06-01 at 23:40 said:}

Ha! In any case, thanks for the pointer to the wheel! This ass continues to meander to one bale or the other! I must enjoy watching him do it.

\end{sffamily}
\end{footnotesize}

